\scanpage{95}

S'il y a aussi la structure motivique de Tate, définissant donc

\[
\varepsilon:G'\unsure{-|\longrightarrow}G_m
\]

alors $\varepsilon i_1(\lambda)=\lambda^n$, et l'entier $n$ est le degré du \deleted{$\unsure{1(-1)}$} ``motif de Tate ``effectif'' (\unsure{moralement|vraisblmt|réalisant|noté} $T_{\ell}(G_m)^{-1}$). Dans le cas auquel nous nous intéressons, cela doit être $1$ (on dit qu'il y a compatibilité avec la structure de Tate ---) :

\[
\varepsilon i_1(\lambda)=\lambda
\]

\uline{Exemple} Soit $S$ préschéma \inserted{connexe} sur $\mathbb Q$, \deleted{\unsure{l}} et $\mathcal M$ la catégorie des motifs sur $S$. Le \uline{foncteur de} \uline{De Rham} sur $\mathcal M$ est alors à valeurs dans les \deleted{\unsure{fais|f}}Modules loc. libres sur $S$, et est admissible.

À ce titre, il définit donc des données $F'$ (Ce n'est autre que le \uline{foncteur de Hodge}) \deleted{\unsure{défini}} $G'$, $i:G_m\longrightarrow G'$, (d'où $\unsure{G'^{>}|B'^{>}|\mathfrak g'^{>}}$) et un $G'^{>}$-torseur $Q$.

\unsure{Etant|Étant} donné la structure de $G'^{>}$ (groupe unipotent), il s'ensuit d'ailleurs que si \uline{$S$ est réduit à un pt}, $Q$ \uline{est trivial}, donc le foncteur de de-Rham filtré \uline{splitte} (mais le splitting du foncteur n'est pas canonique, en particulier n'est pas stable par les autom de ce foncteur, \unsure{ni|même|ou} par ceux qui respectent la filtration
% archive sheet number: 94
